\customchapter{RESUMEN}

En Lima Metropolitana existen puntos críticos de congestión vehicular, especialmente en intersecciones y corredores urbanos de alto tráfico. Esta investigación propone \textbf{reducir el retraso promedio por vehículo} mediante \textbf{semaforización adaptativa} con \textbf{Aprendizaje por Refuerzo (RL)} y validación en \textbf{SUMO}.

La solución consiste en un \textit{pipeline} reproducible con cuatro unidades: (U1) datos y preprocesamiento; (U2) agentes de \textbf{RL} (un agente por intersección); (U3) simulación y medición en SUMO; y (U4) \textbf{coordinación distribuida entre semáforos}. En esta fase se implementaron U1, U2 y U3 y se validaron sobre \textbf{cuatro escenarios sintéticos} construidos con \texttt{netconvert}: un cruce de dos fases, un cruce con giros protegidos, un corredor de dos cruces y una red de tres cruces con rotonda. El algoritmo implementado es \textbf{Q-learning} tabular, comparado sobre las mismas diez semillas contra dos controles de referencia: tiempo fijo tuneado y actuado nativo de SUMO. Quedan previstos SARSA, DQN y PPO, el cruce real de Lima importado de OpenStreetMap y el módulo U4, en el que cada agente compartirá mensajes ligeros con sus \textit{vecinos} (fase actual, colas previstas, salida esperada) para \textbf{coordinar ondas verdes} y \textbf{evitar spillback}, con variantes multiagente (IPPO/MAPPO), \textbf{estado extendido} y \textbf{recompensa con término de coordinación}.

La métrica primaria es el \textbf{retraso por vehículo} (\texttt{timeLoss} de SUMO); como secundarias se miden la espera detenida, las colas promedio y máxima, las paradas por vehículo, el \textit{throughput}, las emisiones (CO\textsubscript{2}, NOx) y la recompensa por decisión. Las metas del trabajo completo son una reducción del retraso del \textbf{25–40\%} y de las colas de al menos un \textbf{20\%} frente al control de tiempo fijo. Con Q-learning tabular no se alcanzan: el agente reduce el retraso un 6.4\% en el cruce de dos fases y lo aumenta entre un 6 y un 17\% en los otros tres escenarios, mientras que el control actuado lo reduce entre un 4 y un 34\%. El techo de la representación tabular al crecer el número de fases y la ventaja de un programa fijo coordinado sobre agentes que no se comunican son los resultados que orientan las siguientes fases hacia \textit{Deep RL} y coordinación distribuida.

\noindent \textbf{Palabras clave:}\\
\noindent semáforos inteligentes; aprendizaje por refuerzo; coordinación distribuida; SUMO; Lima Metropolitana; control de tráfico
